\chapter{Motivation and Problem Formulation}
\label{chap:motivation}

\section{Problem Definition}

Let $\mathcal{D}=\{(x_i,y_i)\}_{i=1}^{N}$ be a labeled downstream dataset and let
$f_{\theta_0}$ be a pretrained model. The base parameters $\theta_0$ are frozen. For a set of
candidate linear modules $\mathcal{M}$, an adaptation policy selects
\begin{equation}
  \pi(\mathcal{D}) =
  \left(r_{\mathrm{target}}, r_{\mathrm{init}},
  \alpha, p_{\mathrm{drop}}, \mathcal{M}_{\mathrm{target}}, \mathcal{S}\right),
  \label{eq:policy-output}
\end{equation}
where $\mathcal{S}$ denotes the AdaLoRA schedule.

The ideal policy would maximize validation performance subject to parameter, memory, and runtime
constraints. Directly solving that bilevel optimization would require many training runs.
Star-LoRA instead defines an inexpensive deterministic approximation
\begin{equation}
  \pi(\mathcal{D}) = \phi(z(\mathcal{D})),
\end{equation}
where $z(\mathcal{D})$ is a vector of dataset statistics and $\phi$ is a bounded rule.

\section{Design Hypotheses}

\subsection{Dataset Size and Capacity}

As $N$ increases, the learner observes more task evidence and can support more trainable capacity.
The relationship should not be linear because increasing a dataset from 100 to 1,000 examples is
usually more consequential than increasing it from 10,100 to 11,000. A logarithmic size factor is
therefore used.

\subsection{Distribution Quality}

The raw sample count can be misleading. If one class dominates, many examples provide redundant
label evidence. If a large fraction of tokens occurs only once, token-level estimates are sparse.
If token frequencies are extremely skewed, frequent lexical patterns can dominate learning.
Star-LoRA treats these properties as reasons to reduce rank and increase regularization.

\subsection{Complexity}

Normalized label entropy is high when classes are represented more evenly. For a binary dataset,
this indicates that the model must meaningfully distinguish both classes rather than exploit a
majority label. Input-embedding variance provides a rough model-dependent estimate of diversity in
the input representations. Star-LoRA treats higher complexity as a reason to increase rank.

\subsection{Adapter Surface}

Rank controls capacity inside each selected matrix. The set of selected matrices controls where
the model can adapt. Query and value projections are common LoRA targets. For complex or sparse
data, Star-LoRA also permits key and output projections. For sufficiently large data with strong
token-frequency skew, it additionally permits the gate, up, and down projections of the MLP.
This policy increases expressive power but can be expensive.

\section{Why a Pre-Training Policy?}

A pre-training policy has three possible advantages. First, it provides task-conditioned defaults
without a rank grid search. Second, it can regularize a very small dataset before noisy gradients
are observed. Third, it produces an explicit explanation: a selected rank can be traced to
measured statistics.

It also has clear limitations. Dataset statistics are proxies, thresholds are handcrafted, and
they do not measure actual validation performance. Two tasks with similar label and token
statistics may require different model changes. The policy should therefore be evaluated as a
useful prior, not as an oracle.

\section{Stability as a Separate Question}

The term ``stability-aware'' can refer to two different systems:
\begin{enumerate}
  \item a monitor that computes a stable importance score for analysis; or
  \item an allocator that actually uses that score to make pruning decisions.
\end{enumerate}
The current repository implements the first system. PEFT's AdaLoRA implementation still performs
the actual rank allocation. This distinction prevents the thesis from attributing performance
improvements to an allocator that has not yet been implemented.
